% !TEX TS-program = pdflatex
\documentclass[11pt,a4paper]{article}

\usepackage[margin=2.5cm]{geometry}
\usepackage[T1]{fontenc}
\usepackage{lmodern}
\usepackage{microtype}
\usepackage{setspace}
\onehalfspacing

\usepackage{amsmath,amssymb}
\usepackage{booktabs}
\usepackage{array}
\usepackage{enumitem}
\usepackage[hidelinks]{hyperref}
\usepackage{xurl}
\makeatletter
\g@addto@macro{\UrlBreaks}{\do\/\do-\do_\do.\do=\do?\do&}
\makeatother

\usepackage{titlesec}
\titleformat{\section}{\large\bfseries}{\thesection.}{0.5em}{}
\titleformat{\subsection}{\normalsize\bfseries}{\thesubsection}{0.5em}{}

\title{\textbf{Case Study, AWH Comparison, and Project Scope}\\
\large Atmospheric Water Harvesting for Autonomous PV Panels Cleaning ---\\
Kom Ombo/Benban-Class Utility PV in Upper Egypt}
\author{Graduation Project --- Supporting Notes}
\date{\today}

\begin{document}
\maketitle

\begin{abstract}
\noindent
This note supports \textit{Atmospheric Water Harvesting for Autonomous PV Panels Cleaning} and answers four questions: (1)~Is a $\sim$200--500\,MW Egyptian PV site ``large''? (2)~Which AWH approach fits Aswan/Upper Egypt weather, desert dust, and cleaning-water duty for \textbf{autonomous} cleaning? (3)~How can you compare thermoelectric cooling (TEC) against other AWH families using the same numbers? (4)~Must you design the full robot if the focus is AWH? Worked comparisons use \textbf{illustrative} literature ranges for specific energy consumption (SEC); replace them with your own measured or simulated SEC for the final defense.
\end{abstract}

\tableofcontents
\newpage

\section{Quick answers (read this first)}
\begin{description}[style=nextline]
  \item[Is $\sim$200--500\,MW ``large''?] \textbf{Yes.} In industry terms, \textbf{utility-scale} often starts around tens of MW; \textbf{large utility} commonly means $\geq 50$\,MW single-site or a GW-cluster sub-block. A $\sim$200\,MW or $\sim$500\,MW plant is unambiguously large.
  \item[Best harvest method here?] For \textbf{Upper Egypt (hot, often low daytime RH, seasonal dust)}, the most \textbf{scalable} options are usually \textbf{active dehumidification / vapor-compression-style condensation} (possibly hybridized with storage) or \textbf{sorption/nighttime cycling} if you can integrate heat properly. \textbf{Fog nets} are a poor primary match (fog is rare). \textbf{Passive dew} is usually \textbf{water-thin} per square meter. \textbf{TEC can work as a lab/demo} but is often \textbf{energy-weak} at the same water output vs compressor systems; treat professor preference as a \textbf{benchmark}, not an automatic optimum, unless you prove it with \textbf{your} boundary conditions.
  \item[Do you have to design robots?] \textbf{Not automatically.} \textbf{Autonomous cleaning} can use commercial robots or automated skids; many theses \textbf{define the cleaning interface} (flow rate, pressure, water quality, schedule) and treat the cleaner as a \textbf{black box}. You \textbf{do} need a clear \textbf{water balance} and \textbf{delivery concept}. Full mechatronics design is only required if \textbf{your advisor-approved scope} says so.
\end{description}

\section{Pick a big Egyptian PV anchor (case study)}
\subsection{Recommended primary anchor}
\textbf{Kom Ombo PV complex (Aswan Governorate)} is a strong thesis anchor because it is clearly \textbf{utility-scale}, well publicized, and located in the same broad \textbf{arid soiling} regime as Benban.

\subsection{Useful public links}
\begin{itemize}[leftmargin=*]
  \item ACWA Power project page: \url{https://acwapower.com/en/projects/kom-ombo-pv/}
  \item Benban cluster context (statistics vary by year; verify before quoting exact MW): \url{https://en.wikipedia.org/wiki/Benban_Solar_Park}
  \item Weather / solar resource (site coordinates): \url{https://power.larc.nasa.gov/}
  \item Hourly humidity, wind, rain (grid cell): \url{https://cds.climate.copernicus.eu/}
  \item Quick PV resource benchmarks: \url{https://globalsolaratlas.info}
\end{itemize}

\subsection{What ``dirt'' means here (so you size the right water)}
Upper Egypt PV soiling is dominated by \textbf{mineral dust}, often with \textbf{small particle sizes} that adhere after dew/transient moisture events. Bird droppings and agricultural drift can appear locally. Implications:
\begin{itemize}[nosep]
  \item \textbf{Wet cleaning} is common for full recovery, but \textbf{water chemistry} matters (TDS/hardness $\rightarrow$ spotting).
  \item \textbf{Rain} is infrequent compared to Europe; do not assume rain washes modules clean monthly without evidence.
\end{itemize}

\section{How much water do you need? (repeatable calculation)}
\subsection{Core formulas}
Let $A_{\mathrm{pv}}$ be the total module glass area you clean [m$^2$] (approximate from plant DC power if needed):
\begin{equation}
A_{\mathrm{pv}} \approx \frac{P_{\mathrm{dc,STC}}}{\eta_{\mathrm{mod}}\,1000}
\quad\text{with }P_{\mathrm{dc,STC}}\text{ in W and }\eta_{\mathrm{mod}}\text{ as a fraction.}
\end{equation}
Let $w$ be process water intensity in \textbf{liters per square meter per cleaning event}, and $N_{\mathrm{clean}}$ the number of such events per year. With safety margin $f_{\mathrm{margin}}$:
\begin{equation}
V_{\mathrm{year}} = w\,A_{\mathrm{pv}}\,N_{\mathrm{clean}}\,(1+f_{\mathrm{margin}}),
\qquad
\dot m_{\mathrm{avg}} = \frac{\rho\,V_{\mathrm{year}}}{86400\times 365}
\end{equation}
where $\rho\approx 1000$\,kg/m$^3$ for fresh water and $\dot m_{\mathrm{avg}}$ is a \textbf{calendar-averaged} mass flow [kg/s] (useful for comparing AWH averages).

\textbf{Literature anchor for $w$ (read before quoting):}\\
\url{https://www.sciencedirect.com/science/article/abs/pii/S0038092X20302450}

\subsection{Illustrative numbers (replace with your plant + O\&M interview)}
Assume a \textbf{representative block}: $P_{\mathrm{dc,STC}}=200$\,MW and $\eta_{\mathrm{mod}}=0.215$:
\begin{equation}
A_{\mathrm{pv}}\approx \frac{200\times 10^6}{0.215\times 1000}
\approx 9.30\times 10^5\ \mathrm{m^2}.
\end{equation}
Take \textbf{example} $w=1.8$\,L/m$^2$ per full wet pass, $N_{\mathrm{clean}}=18$/yr, $f_{\mathrm{margin}}=0.25$:
\begin{equation}
V_{\mathrm{year}} = 0.0018 \times (9.30\times 10^5)\times 18 \times 1.25
\approx 3.77\times 10^4\ \mathrm{m^3/yr}.
\end{equation}
Average daily volume: $V_{\mathrm{year}}/365 \approx 103.2$\,m$^3$/day ($\approx 103{,}200$\,kg/day).

\textbf{Important thesis logic:} real plants may use \textbf{partial wetting}, \textbf{brush assist}, or \textbf{recycle loops}, which can cut $w$. Your job is to \textbf{justify} $w$ and $N_{\mathrm{clean}}$ using Egypt dust literature + (ideally) operator data.

\section{Compare AWH methods fairly (same water mass, same SEC definition)}
\subsection{Definition (electric SEC)}
For any electric AWH subsystem producing distilled-like water mass $m_w$ over a period using electricity $E_{\mathrm{el}}$:
\begin{equation}
\mathrm{SEC}_{\mathrm{el}}=\frac{E_{\mathrm{el}}}{m_w}\quad[\mathrm{kWh/kg}]
\quad\Bigl(=\mathrm{kWh/L}\text{ when }\rho=1000\,\mathrm{kg/m^3}\Bigr).
\end{equation}

\subsection{Literature ``typical band'' anchors (use as starting hypotheses)}
A recent synthesis comparing averages across reported studies (your library PDF may differ slightly from these rounded figures) gives \textbf{order-of-magnitude} central tendencies:
\begin{center}
\small
\begin{tabular}{@{}p{4.6cm}cp{4.8cm}@{}}
\toprule
\textbf{Family (active)} & \textbf{SEC$_{\mathrm{el}}$} & \textbf{Source entry point} \\
\midrule
Absorption-style (reported average) & $\approx 0.38$\,kWh/kg &
\href{https://link.springer.com/article/10.1007/s42452-024-05746-z}{Springer (2024)}\\
Cooling condensation / dehumidification (avg.) & $\approx 0.42$\,kWh/kg & same \\
Adsorption-style (avg.) & $\approx 1.16$\,kWh/kg & same \\
\bottomrule
\end{tabular}
\end{center}
\normalsize
\noindent
\textbf{Critical comparability paper (read their metrics before claiming ``best''):}\\
\url{https://pubs.rsc.org/en/content/articlelanding/2022/ee/d2ee03079a}

\subsection{Thermoelectric cooler (TEC) benchmark (professor's option)}
TEC systems can produce condensate by cooling a surface below the dew point, but their \textbf{coefficient of performance} for large $\Delta T$ across stages is typically modest compared to vapor-compression chillers. In student projects, it is acceptable to treat TEC as a \textbf{device family} whose SEC$_{\mathrm{el}}$ you will measure in the prototype band \textbf{0.5--15\,kWh/kg} depending on inlet air state, heat sink design, and whether you count \textbf{only TEC input} vs \textbf{TEC + fans + pumps}. \textbf{Do not invent a single universal SEC for TEC}; pick a conservative design-band, then replace with measurement.

\subsection{Worked comparison using \textbf{one} average daily water target}
Use the illustrative average from Section~3.2: $\dot m_{\mathrm{avg}}\approx 103{,}200$\,kg/day.

Daily electric energy if the AWH block must supply that day's average entirely by electricity at a fixed SEC:
\begin{equation}
E_{\mathrm{day}} = \mathrm{SEC}_{\mathrm{el}}\times m_{\mathrm{day}},
\qquad m_{\mathrm{day}}=\rho\,(V_{\mathrm{year}}/365).
\end{equation}

\begin{center}
\small
\begin{tabular}{@{}lccc@{}}
\toprule
\textbf{Option} & \textbf{SEC$_{\mathrm{el}}$} & \textbf{$m_{\mathrm{day}}$} & \textbf{$E_{\mathrm{day}}$} \\
 & \textbf{[kWh/kg]} & \textbf{[kg]} & \textbf{[MWh]} \\
\midrule
Cooling condensation (illustr.) & 0.42 & 103{,}200 & 43.4 \\
Absorption (illustr.) & 0.38 & 103{,}200 & 39.2 \\
Adsorption (illustr.) & 1.16 & 103{,}200 & 119.7 \\
TEC (conservative band example) & 6.0 & 103{,}200 & 619.2 \\
TEC (pessimistic band example) & 12.0 & 103{,}200 & 1238.4 \\
\bottomrule
\end{tabular}
\end{center}
\normalsize

\noindent\textbf{Interpretation (what to say in your report):} With the \textbf{same} water obligation, TEC often loses on \textbf{kWh per kg} unless your measured prototype (with excellent heat rejection and night-time operation) pushes SEC$_{\mathrm{el}}$ down dramatically. Even then, compare \textbf{peak power} and \textbf{capex} (modules per liter) honestly.

\subsection{Methods that are \textbf{not} mainly ``SEC$_{\mathrm{el}}$'' stories}
\begin{description}[style=nextline,leftmargin=0pt,labelwidth=!]
  \item[Passive dew / radiative collectors] Often reported in small daily yields per m$^2$ of collector (high area footprint). Compare using $Y_w$ [kg/m$^2$/day] from experiments matched to your night RH/temperature.
  \item[Fog collection] Generally \textbf{not} a primary Upper Egypt strategy (fog frequency low). Useful mostly as a ``wrong-tool'' baseline with citation from climatology for your coordinates.
  \item[Sorption with heat integration] If you have excess \textbf{thermal} energy (not only electricity), add a \textbf{thermal SEC} term or a \textbf{primary-energy} metric; otherwise you mis-rank the system.
\end{description}

\section{PV simulation reminder (what PVsyst does / does not do)}
PVsyst models soiling as \textbf{monthly loss factors} you justify; it does not replace an AWH model. Official help:\\
\url{https://www.pvsyst.com/help/project-design/array-and-system-losses/soiling-loss.html}

\section{Project scope: do you design robots?}
\subsection{A clean scope split (recommended)}
\begin{table}[h]
\centering
\begin{tabular}{@{}p{4.2cm}p{10.5cm}@{}}
\toprule
\textbf{In-scope (AWH-centric)} & Water demand model; storage; treatment (filtration/RO/demineralization as needed); energy scheduling; techno-economic comparison; prototype of \textbf{AWH core}. \\
\midrule
\textbf{Optional (only if advisor agrees)} & End-effector robotics, fleet routing, autonomous navigation, structured mechatronics thesis. \\
\midrule
\textbf{Minimum viable interface} & Specify \textbf{peak flow}, \textbf{pressure}, \textbf{water quality limits}, and \textbf{connection point} (tank $\rightarrow$ existing cleaning skid). \\
\bottomrule
\end{tabular}
\end{table}

\subsection{Suggested sentence for your proposal}
\textit{``This work sizes and prototypes atmospheric water harvesting and water treatment to meet a defined PV cleaning water specification; the mechanical cleaning implement is treated as an existing O\&M process unless explicitly expanded in Phase~2.''}

\section{Design Translation and Prototype Package}
\subsection{Why the calculations are not the end goal}
The analytical results in this note exist to identify the \textbf{best candidate configuration} and define a \textbf{traceable engineering basis} for simulation, prototyping, and delivery. The project deliverable is therefore the \textbf{selected buildable system} and its validation workflow, not the comparison table by itself.

\subsection{Calculated best-option parameters and selected build path}
\begin{itemize}[leftmargin=*]
  \item \textbf{Lowest numerical LCOW route:} \textbf{Absorption} using the literature-average SEC band of about 0.38\,kWh/kg.
  \item \textbf{Prototype-selected route:} \textbf{modular vapor-compression condensation with polishing treatment} because it stays close to the best LCOW result while being easier to source, assemble, instrument, and test.
  \item \textbf{Full-scale reference basis:} $P_{\mathrm{dc}}=200$\,MW, $\eta_{\mathrm{mod}}=0.215$, $A_{\mathrm{pv}}\approx 9.30\times 10^5$\,m$^2$, $w=1.8$\,L/m$^2$/event, $N_{\mathrm{clean}}=18$/yr, $f_{\mathrm{margin}}=0.25$, and $V_{\mathrm{year}}\approx 3.77\times 10^4$\,m$^3$/yr.
  \item \textbf{Scaled prototype target:} 200\,L/day condensate, treatment polishing, 500\,L raw tank, 500\,L polished-water tank, and a 2.0\,m$^3$/h delivery pump at 3.0\,bar for interface validation.
\end{itemize}

\subsection{Traceability matrix}
\begin{table}[h]
\centering
\small
\begin{tabular}{@{}p{4.0cm}p{4.0cm}p{6.8cm}@{}}
\toprule
\textbf{Calculation} & \textbf{Decision} & \textbf{Traceable design outcome} \\
\midrule
$V_{\mathrm{year}}\approx 3.77\times 10^4$\,m$^3$/yr & Use modular generation plus storage & Requires buffered production instead of direct one-pass supply only. \\
$m_{\mathrm{day,avg}}\approx 103.2$\,m$^3$/day & Define prototype scaling ratio & Sets the 200\,L/day prototype validation target. \\
Absorption SEC/LCOW minimum & Keep as benchmark & Serves as the analytic performance floor for later design improvements. \\
VCR remains close to optimum & Choose build route & Supports an immediately buildable prototype with standard components. \\
Condensate quality trend is favorable & Add polishing, not full desalination by default & Fixes the treatment train to sediment + carbon + UV with optional RO bypass. \\
\bottomrule
\end{tabular}
\end{table}

\subsection{Simulation plan}
\begin{enumerate}[leftmargin=*]
  \item Assemble an hourly weather dataset for Benban from NASA POWER and ERA5.
  \item Validate the cleaning schedule in PVsyst using the same soiling assumptions used to derive the water duty.
  \item Simulate the VCR prototype with compressor cycling, fan power, condensate collection, storage-tank level, and polishing pressure drop.
  \item Run summer peak, low-RH night, filter-fouling, and missed-cleaning-cycle recovery scenarios.
  \item Approve the design only if 95\% of simulated days meet the condensate target and final-water TDS remains within the rinse specification.
\end{enumerate}

\subsection{Engineering drawings, bill of materials, and component specifications}
\textbf{Drawing set required for prototype release:}
\begin{itemize}[leftmargin=*]
  \item D-101 process flow diagram.
  \item D-102 skid general arrangement (2200\,mm $\times$ 1400\,mm $\times$ 2100\,mm).
  \item D-103 condensate tank and drain routing.
  \item D-104 treatment manifold and sampling-point layout.
  \item D-105 electrical one-line and control panel layout.
  \item D-106 instrumentation loop diagram.
  \item D-107 cleaning-water interface manifold for the autonomous cleaning test rig.
\end{itemize}

\textbf{BOM and specification targets:}
\begin{itemize}[leftmargin=*]
  \item Hermetic compressor package sized for about 200\,L/day condensate.
  \item Hydrophilic evaporator coil, air-cooled condenser, and two EC axial fans.
  \item 500\,L raw-water tank, 500\,L polished-water tank, and stainless or food-safe wetted materials.
  \item 5\,$\mu$m sediment filter, activated-carbon cartridge, UV disinfection unit, and optional RO bypass spool.
  \item Delivery pump rated 2.0\,m$^3$/h at 3.0\,bar, with flow meter, pressure gauge, inline TDS, T/RH sensors, level switches, and PLC or datalogger.
\end{itemize}

\subsection{Build manual, safety, and quality control}
\begin{enumerate}[leftmargin=*]
  \item Assemble and leak-test the refrigeration loop before electrical energization.
  \item Install treatment, storage, and pump hardware in the sequence defined on D-101.
  \item Calibrate sensors and execute dry commissioning, wet commissioning, and a 24-hour condensate trial.
  \item Enforce lockout-tagout, refrigerant handling, electrical earthing, splash protection, and water-tank sanitation protocols.
  \item Hold mandatory QC checkpoints at leak test, airflow verification, condensate-rate verification, TDS verification, and discharge-pressure verification.
\end{enumerate}

\subsection{Updated master project plan}
\begin{table}[h]
\centering
\small
\begin{tabular}{@{}p{2.0cm}p{4.2cm}p{3.1cm}p{2.0cm}p{3.4cm}@{}}
\toprule
\textbf{Milestone} & \textbf{Deliverable} & \textbf{Resources} & \textbf{Date} & \textbf{Success criteria} \\
\midrule
M1 & Freeze inputs, route ranking, and traceability matrix & PI + analyst & T0+2d & Design basis approved \\
M2 & Complete PVsyst and hourly AWH simulations & PV analyst + controls student & T0+7d & Weather-linked model closed \\
M3 & Release drawing package D-101 to D-107 & CAD lead + mechanical lead & T0+10d & Interfaces fully dimensioned \\
M4 & Procure BOM and fabricate prototype skid & Procurement + workshop technician & T0+17d & Long-lead items inspected \\
M5 & Build and commission prototype & Mechanical + electrical + safety officer & T0+24d & Stable 24-hour safe run \\
M6 & Validate against calculated targets & Test engineer + QA lead & T0+28d & Production, TDS, and pressure targets achieved \\
\bottomrule
\end{tabular}
\end{table}

\section{Bibliography pointers (export properly in Bib\TeX)}
\begin{itemize}[leftmargin=*]
  \item PV soiling / cleaning optimization entry point: \url{https://www.sciencedirect.com/science/article/abs/pii/S0306261921000039}
  \item AWH metrics: \url{https://pubs.rsc.org/en/content/articlelanding/2022/ee/d2ee03079a}
  \item SEC averages (absorption/condensation/adsorption): \url{https://link.springer.com/article/10.1007/s42452-024-05746-z}
  \item Psychrometrics: ASHRAE Fundamentals Handbook (Chapter~1).
\end{itemize}

\end{document}
